%==============================================================================
% Sjabloon poster bachproef
%==============================================================================

\documentclass[a0,portrait]{hogent-poster}

\course{Bachelorproef}
\studyprogramme{toegepaste informatica}
\academicyear{2025-2026}
\institution{Hogeschool Gent, Valentin Vaerwyckweg 1, 9000 Gent}

\title{Optimalisatie en uniformering van documentatie binnen Modern Workplace Infrastructure bij ESC bv}
\subtitle{Geautomatiseerde migratie van Microsoft OneNote naar Remote Desktop Manager}
\author{Wim Meirlaen}
\email{wim.meirlaen@student.hogent.be}
\supervisor{Dhr. A. Veldeman}
\cosupervisor{Dhr. K. De Jaeger (E.S.C. BV)}

\specialisation{Systeem- en Netwerkbeheer}
\keywords{OneNote, Remote Desktop Manager, Microsoft Graph API, PowerShell, migratie, PAM}
\projectrepo{https://github.com/WimMeirlaen/BP}

\begin{document}

\maketitle

\begin{abstract}
E.S.C.\ BV beheert klantdocumentatie voor honderden klanten in Microsoft OneNote, verspreid over meerdere SharePoint-locaties per klant. Dit leidt tot zoekproblemen, ontbrekende toegangscontrole en veiligheidsrisico's. Na de overstap naar Remote Desktop Manager (RDM) werd een geautomatiseerde PowerShell-migratiepipeline ontwikkeld via Microsoft Graph API, gevalideerd voor meerdere klantvaults in productie.
\end{abstract}

\begin{multicols}{2}

\section{Probleemstelling}

Klantdocumentatie bij E.S.C.\ BV is verspreid over \textbf{drie SharePoint-locaties} per klant:

\begin{itemize}
  \item Medewerkers moeten op drie plaatsen zoeken zonder garantie alles te vinden
  \item Credentials staan als \textbf{platte tekst} — geen versleuteling of logging per item
  \item \textbf{Least-privilege} niet haalbaar: toegang tot een site = toegang tot alle inhoud
\end{itemize}

\textbf{Remote Desktop Manager (RDM)} biedt vaults per klant met toegangscontrole en auditlogging, in lijn met PAM-principes.

\section{Onderzoeksvraag}

Hoe kan klantinformatie \textbf{geautomatiseerd, veilig en schaalbaar} worden gemigreerd van Microsoft OneNote naar Remote Desktop Manager, rekening houdend met API-beperkingen en beveiligingsvereisten?

\section{Methode}

Drie extractiemethoden werden geëvalueerd. Microsoft Graph API werd gekozen wegens betrouwbaarheid en schaalbaarheid; COM-automatisering bleek geblokkeerd door beveiligingsbeleid; handmatige export is niet haalbaar bij honderden klanten.

Authenticatie via \textbf{Entra~ID} met device code flow en gedelegeerde permissies — geen clientgeheim in het script.

\section{Architectuur van de pipeline}

\begin{center}
\begin{minipage}{\linewidth}
  \centering
  \includegraphics[width=1.0\linewidth]{fig-06-two-phase-architectuur}
  \captionsetup{type=figure}
  \captionof{figure}{Two-phase architectuur: exportketen (links) en importketen (rechts)}
\end{minipage}
\end{center}

\begin{itemize}
  \item \textbf{Export}: haalt tekst, afbeeldingen en bijlagen op via Graph API
  \item \textbf{Deduplicatie}: per sectie/paginatitel enkel de meest recente versie bewaard
  \item \textbf{Throttlingbeheer}: circuit breaker en exponentiële back-off
  \item \textbf{Import}: plaatst HTML-bestanden in de juiste klantvault via Devolutions.PowerShell; structuur en inhoud blijven doorzoekbaar in RDM
\end{itemize}

\section{Resultaten}

\begin{itemize}
  \item Volledige export van tekst, afbeeldingen en bijlagen over meerdere SharePoint-sites
  \item Stabiele batchverwerking van meerdere klantvaults \textbf{zonder gebruikersinput}
  \item Handmatig inlogproces (Teams-lidmaatschap, sites zoeken) volledig geautomatiseerd
\end{itemize}

\begin{center}
\begin{minipage}{\linewidth}
  \centering
  \includegraphics[width=0.90\linewidth]{boomstructuur}
  \captionsetup{type=figure}
  \captionof{figure}{OneNote-hiërarchie (secties, pagina's) herkenbaar terug in RDM}
\end{minipage}
\end{center}

\section{Conclusie}

Geautomatiseerde migratie van OneNote naar RDM is \textbf{technisch haalbaar}. Deduplicatie, throttlingbeheer en structuurbehoud maken het verschil tussen een experimenteel script en een productieoplossing. Het resultaat biedt een directe basis voor bedrijfsbrede uitrol en is relevant voor IT-dienstverleners met gelijkaardige migratienoden.

\textit{Vervolgonderzoek:} fijnmazigere API-permissies, AI-gestuurde categorisatie (lokaal taalmodel).

\end{multicols}
\end{document}
